\paragraph{}Rozpoznawanie tożsamości na podstawie obrazu twarzy (ang. Face Recognition) jest jednym z kluczowych zagadnień współczesnej wizji komputerowej i znajduje szerokie zastosowanie w systemach bezpieczeństwa, monitoringu wizyjnego oraz urządzeniach konsumenckich. W ostatnich latach, dzięki rozwojowi głębokich sieci neuronowych oraz funkcji strat opartych na marginesach kątowych, takich jak ArcFace, nowoczesne modele typu state-of-the-art osiągają bardzo wysoką skuteczność w warunkach kontrolowanych i na standardowych zbiorach testowych.

Jednocześnie skuteczność systemów rozpoznawania twarzy spada w obecności zakłóceń charakterystycznych dla rzeczywistych scenariuszy. Szczególnym wyzwaniem jest częściowa okluzja twarzy, w tym zasłonięcie obszaru oczu spowodowane noszeniem okularów, masek ochronnych lub ograniczoną widocznością w systemach monitoringu. Chociaż złożone obliczeniowo modele typu state-of-the-art wykazują relatywnie wysoką odporność na tego typu zakłócenia, ich lżejsze warianty, projektowane z myślą o zastosowaniach mobilnych, jak również starsze architektury, często charakteryzują się znacznym spadkiem skuteczności identyfikacji w warunkach okluzji.

Celem niniejszej pracy jest analiza wpływu częściowej okluzji twarzy na skuteczność
systemów rozpoznawania tożsamości oraz zaproponowanie modyfikacji architektury
opartej na mechanizmach atencji, zwiększającej odporność modeli na tego typu zakłócenia.